\documentclass{article}
\usepackage[utf8]{inputenc}

\title{Paratexts: Thresholds of Interpretation}
\author{G\'{e}rard Genette}
\date{1997}

\begin{document}
\maketitle

\pgmark{1}\section*{Introduction}

A literary work consists, entirely or essentially, of a text. But this
text is rarely presented in an unadorned state, unreinforced and
unaccompanied by a certain number of verbal or other productions, such as
an author's name, a title, a preface, illustrations.

\pgmark{2}These productions, which I have proposed to call by the
generic term \emph{paratext}, surround and prolong the text precisely in
order to present it.

\pgmark{3}\section{The peritext and the epitext}

The paratext is divided into the \emph{peritext} (everything physically
attached to the text: title, preface, footnotes) and the \emph{epitext}
(everything outside the volume: interviews, letters).

\pgmark{4}This excerpt is a stub for the Virgil Library dev preview ---
enough text to confirm that a second paper renders alongside the first
when selected from the catalog.

\end{document}
